\lecture{27}{Fri 04 Nov 2022 11:31}{}

\subsection{One-sided limits}

Assume $c \in [a,b)$.
We define the \textit{right limit} of $f$ at $c$ as
\[
	\lim_{x \to c+} f(x) = \lim_{x \to c; x \in [c,b]} f(x) = \lim_{x \to c} f|_{[c,b]}
.\] 

Similarly if $c \in (a,b]$, we define
\[
	\lim_{x \to c-} f(x) = \lim_{x \to c; x \in [a,c]} f(x) = \lim_{x \to c} f|_{[a,c]}
.\] 

\begin{proposition}
	If $c \in (a,b)$, then $\lim_{x \to c} f(x)$ exists iff $\lim_{x \to c-} f(x)$ and $\lim_{x \to c+} f(x)$ both exist and are equal to each other.
\end{proposition}
\begin{proof}
	Very straightforward using $\varepsilon-\delta$ definition.
\end{proof}

\begin{example}
	(Monotone Functions)
	We say $f: [a,b] \to \mathbb{R}$ is \textit{nondecreasing} if $x\le y$ implies $f(x) \le f(y)$ for all  $x,y \in E$.

	We say $f$ is \textit{strictly increasing} if $x<y$ implies $f(x) <f(y)$.
\end{example}

\begin{proposition}
	Let $f: [a,b] \to \mathbb{R}$ be increasing on $[a,b]$. Then 
	\begin{enumerate}
		\item 
	for any $c \in [a,b)$ we have $f(c+):= \lim_{x \to c+} f(x)$ exists. Moreover, $f(c+) = \operatorname{inf}_{x \in (c,b]}f(x)$.
\item 
	Similarly, if $c \in (a,b]$, we have $f(c-) = \operatorname{sup}_{x \in [a,c)}f(x)$.
\item For $c \in (a,b)$,
	\[
	f(c-) \le f(c) \le f(c+)
.\] 
\item If $c < d, c \in [a,b), d \in (a,b]$, then
	\[
		f(c+) \le f(d-)
	.\] 
	\end{enumerate}
\end{proposition}
\begin{proof}
	\begin{enumerate}
		\item The set $S = \{f(x): c < x \le b\}$ is bounded below by $f(c)$, hence $\alpha:= \operatorname{inf} S = \operatorname{inf}_{c<x\le b}f(x)$ exists.

			Note: $f(c)\le \alpha$ since $f(c)$ is a lower bound. We now claim $\alpha = \lim_{x \to c+} f(x)$. Fix any $\varepsilon>0$. By definition of inf, there is $x_\varepsilon$ such that $c<x_\varepsilon\le b$ and $\alpha \le f(x_\varepsilon)< d + \varepsilon$. 
			Then if $c < x < x_\varepsilon$, we have $\alpha \le f(x) \le f(x_\varepsilon) < \alpha + \varepsilon$.
			In particular this means  $|f(x) - \alpha| < \varepsilon$. If we define $\delta(\varepsilon) = x_\varepsilon - c$, we have
			\[
				c<x<c+\delta(\varepsilon) \implies |f(x) - \alpha| < \varepsilon
			.\] 
			Hence $\lim_{x \to c+} f = \alpha$.
	\item Symmetric.
	\item Showed in part of proof for $(1)$.
	\item Pick $e \in [c,d]$. Then
		\[
			f(c+) = \operatorname{inf}_{(c,b]} f(x) \le  f(e) \le  \operatorname{sup}_{[a,d)} f = f(d-)
		.\] 
	\end{enumerate}
\end{proof}

\begin{theorem}
	If $f: [a,b] \to \mathbb{R}$ is monotone (nondecreasing or nonincreasing), then $f$ is continuous at all but at most countably many points in  $[a,b]$.
\end{theorem}
\begin{proof}
	Assume $f$ is increasing. If $c \in (a,b)$ is a point of discontinuity, then $f(c-) < f(c+)$. Then there exists a rational $r \in (f(c-), f(c+))$. Let $S$ be the set of discontinuity of $f$. We obtain the injection
	 \begin{align*}
		\varphi: S &\longrightarrow \mathbb{Q} \\
		c &\longmapsto \varphi(c) = r_c
	.\end{align*}

	We know $r_{c_1} \neq r_{c_2}$ whenver $c_1\neq c_2$, since if  $c_1<c_2$, we have $r_{c_1}<f(c_1+) < f(c_2-) < f_{c_2}$.
\end{proof}
\begin{proof}
	(Simpler) Define $J_c = (f(c-), f(c+))$. We verify that $\mathcal{F} = \{J_c\}_{f \text{ discontinuous at }c}$ is a disjoint family of open nonempty intervals in $\mathbb{R}$, hence it is at most countable.
\end{proof}

\subsection{Limsup and Liminf at a point}

Let $f: E \to \mathbb{R}, c$ is a cluster point of $E$.

\begin{definition}
	We say $M$ is a \textit{local upper bound} of $f$ near $c$ if there is an open set $U$ such that \[
		f(x) \le M \impliedby x \in \left( U \setminus \{c\}  \right) \cap E
	.\] 
\end{definition}
Equivalently,
\begin{proposition}
	$M$ is a \textit{local upper bound} of $f$ near $c$ iff there is $r>0$ such that for any $x \in E$,
	 \[
		f(x) \le M \impliedby 0 < |x-c| < r
	.\] 
\end{proposition}

\begin{definition}
	We define
	\[
		\limsup_{x \to c} f =
		\operatorname{inf} \{M: M \text{ local upper bound of $f$ near $c$}\}
	.\] 
\end{definition}
\begin{proposition}
	For $r>0$ define
	\[
		\varphi(r) = \operatorname{sup} \{f(x): x \in \left( B_r(c)\setminus \{c\} \right) \cap E \}
	\]
	Then 
	\[
		\limsup_{x \to c} f(x) = \lim_{r \to 0+} \varphi(r) = \operatorname{inf}_{r>0}\varphi(r)
	.\] 
	So formally,
	\[
		\limsup_{x \to c} f(x) = \lim_{r \to 0+} \operatorname{sup}_{\left( B_r(c) \setminus \{c\} \right) \cap E }f
	.\] 
\end{proposition}

\begin{proposition}
	Let $f,g : E \to \mathbb{R}$, $c$ cluster point of $E$. Then \[
		\limsup_{x \to c} f(x) + g(x) \le 
		\limsup_{x \to c} f(x) + \limsup_{x \to c} g(x)
	.\] 
\end{proposition}
\begin{proof}
	Let $x \in \left( B_r(c) \setminus \{c\}  \right) \cap E =: E_r$. Then
	\[
		f(x) + g(x) \le \operatorname{sup}_{E_r}f + \operatorname{sup}_{E_r} g
	.\] 
	Then
	\[
		\operatorname{sup}_{E_r} f(x) + g(x) \le \operatorname{sup}_{E_r}f + \operatorname{sup}_{E_r} g
	.\]
	Hence
	 \[
		\varphi_{f+g}(r) \le \varphi_f(r) + \varphi_g(r)
	.\] 
	Therefore
	\[
		\lim_{r \to 0+} \varphi_{f+g} (r) \le \lim_{r \to 0+} \varphi_f(r) + \lim_{r \to 0+} \varphi_g(r)
	.\] 
	And the result follows.
\end{proof}
\begin{remark}
	Other results of limsup and liminf can also be generalized. Should be able to prove them myself.
\end{remark}

